\documentclass[10pt,a4paper]{amsart}
\usepackage{amsmath,amssymb,amsthm,graphicx}
\usepackage{verbatim}
\newenvironment{reference}{\comment}{\endcomment}
%\newenvironment{reference}{}{}
\newenvironment{slogan}{\comment}{\endcomment}
\newenvironment{history}{\comment}{\endcomment}

% For commutative diagrams we use Xy-pic
\usepackage[all]{xy}

% We use 2cell for 2-commutative diagrams.
\xyoption{2cell}
% Theorem environments.
%
\theoremstyle{plain}
\newtheorem{theorem}[subsection]{Theorem}
\newtheorem{proposition}[subsection]{Proposition}
\newtheorem{lemma}[subsection]{Lemma}

\theoremstyle{definition}
\newtheorem{definition}[subsection]{Definition}
\newtheorem{example}[subsection]{Example}
\newtheorem{exercise}[subsection]{Exercise}
\newtheorem{situation}[subsection]{Situation}

\theoremstyle{remark}
\newtheorem{remark}[subsection]{Remark}
\newtheorem{remarks}[subsection]{Remarks}

\numberwithin{equation}{subsection}

% Macros
%
\def\lim{\mathop{\mathrm{lim}}\nolimits}
\def\colim{\mathop{\mathrm{colim}}\nolimits}
\def\Spec{\mathop{\mathrm{Spec}}}
\def\Hom{\mathop{\mathrm{Hom}}\nolimits}
\def\Ext{\mathop{\mathrm{Ext}}\nolimits}
\def\SheafHom{\mathop{\mathcal{H}\!\mathit{om}}\nolimits}
\def\SheafExt{\mathop{\mathcal{E}\!\mathit{xt}}\nolimits}
\def\Sch{\mathit{Sch}}
\def\Mor{\mathop{\mathrm{Mor}}\nolimits}
\def\Ob{\mathop{\mathrm{Ob}}\nolimits}
\def\Sh{\mathop{\mathit{Sh}}\nolimits}
\def\NL{\mathop{N\!L}\nolimits}
\def\CH{\mathop{\mathrm{CH}}\nolimits}
\def\proetale{{pro\text{-}\acute{e}tale}}
\def\etale{{\acute{e}tale}}
\def\QCoh{\mathit{QCoh}}
\def\Ker{\mathop{\mathrm{Ker}}}
\def\Im{\mathop{\mathrm{Im}}}
\def\Coker{\mathop{\mathrm{Coker}}}
\def\Coim{\mathop{\mathrm{Coim}}}

% Boxtimes
%
\DeclareMathSymbol{\boxtimes}{\mathbin}{AMSa}{"02}

%
% Macros for moduli stacks/spaces
%
\def\QCohstack{\mathcal{QC}\!\mathit{oh}}
\def\Cohstack{\mathcal{C}\!\mathit{oh}}
\def\Spacesstack{\mathcal{S}\!\mathit{paces}}
\def\Quotfunctor{\mathrm{Quot}}
\def\Hilbfunctor{\mathrm{Hilb}}
\def\Curvesstack{\mathcal{C}\!\mathit{urves}}
\def\Polarizedstack{\mathcal{P}\!\mathit{olarized}}
\def\Complexesstack{\mathcal{C}\!\mathit{omplexes}}
% \Pic is the operator that assigns to X its picard group, usage \Pic(X)
% \Picardstack_{X/B} denotes the Picard stack of X over B
% \Picardfunctor_{X/B} denotes the Picard functor of X over B
\def\Pic{\mathop{\mathrm{Pic}}\nolimits}
\def\Picardstack{\mathcal{P}\!\mathit{ic}}
\def\Picardfunctor{\mathrm{Pic}}
\def\Deformationcategory{\mathcal{D}\!\mathit{ef}}
\usepackage{cleveref}
\crefname{originalreference}{original reference}{original references}
\Crefname{originalreference}{Original reference}{Original references}
\makeatletter
\newlabel{situation-versal}{{0DQU}{source}{}{sourcepointer-1}{}}
\newlabel{situation-versal@cref}{{[originalreference][0][]0DQU}{[1][0][]source}{}{}{}}
\makeatother
\crefname{originalreference}{original reference}{original references}
\Crefname{originalreference}{Original reference}{Original references}
\makeatletter
\newlabel{artin-definition-versal}{{07XF}{source}{}{sourcepointer-2}{}}
\newlabel{artin-definition-versal@cref}{{[originalreference][0][]07XF}{[1][0][]source}{}{}{}}
\makeatother
\crefname{originalreference}{original reference}{original references}
\Crefname{originalreference}{Original reference}{Original references}
\makeatletter
\newlabel{artin-lemma-versality-matches}{{0CXK}{source}{}{sourcepointer-3}{}}
\newlabel{artin-lemma-versality-matches@cref}{{[originalreference][0][]0CXK}{[1][0][]source}{}{}{}}
\makeatother
\crefname{originalreference}{original reference}{original references}
\Crefname{originalreference}{Original reference}{Original references}
\makeatletter
\newlabel{artin-lemma-formally-smooth-on-deformation-categories}{{07WK}{source}{}{sourcepointer-4}{}}
\newlabel{artin-lemma-formally-smooth-on-deformation-categories@cref}{{[originalreference][0][]07WK}{[1][0][]source}{}{}{}}
\makeatother
\crefname{originalreference}{original reference}{original references}
\Crefname{originalreference}{Original reference}{Original references}
\makeatletter
\newlabel{spaces-morphisms-lemma-finite-type-points-surjective-morphism}{{06EJ}{source}{}{sourcepointer-5}{}}
\newlabel{spaces-morphisms-lemma-finite-type-points-surjective-morphism@cref}{{[originalreference][0][]06EJ}{[1][0][]source}{}{}{}}
\makeatother
\crefname{originalreference}{original reference}{original references}
\Crefname{originalreference}{Original reference}{Original references}
\makeatletter
\newlabel{artin-lemma-base-change-versal}{{0CXM}{source}{}{sourcepointer-6}{}}
\newlabel{artin-lemma-base-change-versal@cref}{{[originalreference][0][]0CXM}{[1][0][]source}{}{}{}}
\makeatother
\title{Four equivalent characterizations of versality at a finite type point; the completed local ring is a versal ring}
\author{The Stacks Project Authors}
\date{stacks-project:stacks-geometry.tex:local-snapshot-20261012; GNU FDL 1.2 or later}
\begin{document}
\maketitle
\noindent\textit{Editorial note.} All delivered mathematical text is the author's own, copied verbatim from the pinned Stacks Project chapter file stacks-geometry.tex; no statement or proof step was rewritten, reordered or removed.\par\medskip
\noindent\textit{Editorial note.} Span recomputation by label pairing (LANE\_BRIEF 10.2): in this chapter the author's proof environment follows the statement environment as a standalone block; the spec covers the complete environment pair (statement span, proof span) as recorded in the pool candidate.\par\medskip
\noindent\textit{Editorial note.} Occupation scan: combined claim\_key + statement\_anchor + generic-line scan over all registration files (active and withdrawn) found 0 hits for stacks-geometry.tex (ge-occupied.json, 2026-10-14); 207 route=geometry / chapter-list field hits in focused-seeds are alias noise (route name and alphabetical chapter enumerations), excluded per the ladder alias-blind-spot caveat.\par\medskip
\ifdefined\hypertarget\hypertarget{sourcepointer-1}{}\fi
\noindent\textit{Source reference.} Original reference 0DQU: sources/stacks-chow/tags.txt, source lines 17199--17199. Cross-chapter standard prerequisite cited by the author inside the counted span; it is used as a pointer and its own proof is not part of this item.\par\medskip
\ifdefined\hypertarget\hypertarget{sourcepointer-2}{}\fi
\noindent\textit{Source reference.} Original reference 07XF: sources/stacks-chow/tags.txt, source lines 10075--10075. Cross-chapter standard prerequisite cited by the author inside the counted span; it is used as a pointer and its own proof is not part of this item.\par\medskip
\ifdefined\hypertarget\hypertarget{sourcepointer-3}{}\fi
\noindent\textit{Source reference.} Original reference 0CXK: sources/stacks-chow/tags.txt, source lines 16210--16210. Cross-chapter standard prerequisite cited by the author inside the counted span; it is used as a pointer and its own proof is not part of this item.\par\medskip
\ifdefined\hypertarget\hypertarget{sourcepointer-4}{}\fi
\noindent\textit{Source reference.} Original reference 07WK: sources/stacks-chow/tags.txt, source lines 10045--10045. Cross-chapter standard prerequisite cited by the author inside the counted span; it is used as a pointer and its own proof is not part of this item.\par\medskip
\ifdefined\hypertarget\hypertarget{sourcepointer-5}{}\fi
\noindent\textit{Source reference.} Original reference 06EJ: sources/stacks-chow/tags.txt, source lines 8218--8218. Cross-chapter standard prerequisite cited by the author inside the counted span; it is used as a pointer and its own proof is not part of this item.\par\medskip
\ifdefined\hypertarget\hypertarget{sourcepointer-6}{}\fi
\noindent\textit{Source reference.} Original reference 0CXM: sources/stacks-chow/tags.txt, source lines 16212--16212. Cross-chapter standard prerequisite cited by the author inside the counted span; it is used as a pointer and its own proof is not part of this item.\par\medskip
\begin{lemma}
\label{lemma-compare-versal-ring-completion}
In Situation \ref{situation-versal} let $x : U \to \mathcal{X}$ be a
morphism where $U$ is a scheme locally of finite type over $S$.
Let $u_0 \in U$ be a finite type point.
Set $k = \kappa(u_0)$ and denote $x_0 : \Spec(k) \to \mathcal{X}$
the induced map. The following are equivalent
\begin{enumerate}
\item $x$ is versal at $u_0$
(Artin's Axioms, Definition \ref{artin-definition-versal}),
\item $\hat x : \mathcal{F}_{U, k, u_0} \to \mathcal{F}_{\mathcal{X}, k, x_0}$
is smooth,
\item the formal object associated to
$x|_{\Spec(\mathcal{O}_{U, u_0}^\wedge)}$ is versal, and
\item there is an open neighbourhood $U' \subset U$ of $x$ such that
$x|_{U'} : U' \to \mathcal{X}$ is smooth.
\end{enumerate}
Moreover, in this case the completion $\mathcal{O}_{U, u_0}^\wedge$
is a versal ring to $\mathcal{X}$ at $x_0$.
\end{lemma}

\begin{proof}
Since $U \to S$ is locally of finite type (as a composition of such morphisms),
we see that $\Spec(k) \to S$ is of finite type (again as a composition).
Thus the statement makes sense. The equivalence of (1) and (2)
is the definition of $x$ being versal at $u_0$.
The equivalence of (1) and (3) is
Artin's Axioms, Lemma \ref{artin-lemma-versality-matches}.
Thus (1), (2), and (3) are equivalent.

\medskip\noindent
If $x|_{U'}$ is smooth, then the functor
$\hat x : \mathcal{F}_{U, k, u_0} \to \mathcal{F}_{\mathcal{X}, k, x_0}$
is smooth by Artin's Axioms, Lemma
\ref{artin-lemma-formally-smooth-on-deformation-categories}.
Thus (4) implies (1), (2), and (3).
For the converse, assume $x$ is versal at $u_0$.
Choose a surjective smooth morphism $y : V \to \mathcal{X}$ where $V$
is a scheme. Set $Z = V \times_\mathcal{X} U$ and pick a finite type
point $z_0 \in |Z|$ lying over $u_0$ (this is possible by
Morphisms of Spaces, Lemma
\ref{spaces-morphisms-lemma-finite-type-points-surjective-morphism}).
By Artin's Axioms, Lemma \ref{artin-lemma-base-change-versal}
the morphism $Z \to V$ is smooth at $z_0$.
By definition we can find an open neighbourhood $W \subset Z$
of $z_0$ such that $W \to V$ is smooth. Since $Z \to U$ is open,
let $U' \subset U$ be the image of $W$. Then we see that
$U' \to \mathcal{X}$ is smooth by our definition of smooth morphisms
of stacks.

\medskip\noindent
The final statement follows from the definitions as
$\mathcal{O}_{U, u_0}^\wedge$
prorepresents $\mathcal{F}_{U, k, u_0}$.
\end{proof}

\end{document}
